\section{Guía de lectura: entrevistar a un CEO como si fuera un gran modelo MoE}

Esta sección establece primero el método de lectura de toda la entrevista. En la apertura, Zhang Xiaojun (张小珺) propone un marco muy útil: tratar a Li Xiang (李想) como un «gran modelo CEO» y suponer que tiene una arquitectura MoE (Mixture of Experts, mezcla de expertos). Las tres primeras rondas invocan, en orden, al experto técnico, al experto en estrategia y al experto en organización; en la segunda mitad, las preguntas avanzan del modelo, el automóvil y los robots hacia la energía, las relaciones íntimas, la memoria y la sabiduría. Este marco no es una broma: convierte una entrevista de tres horas sobre una persona en una reflexión sistemática sobre las empresas, los productos y las personas en la era de la IA.

El objetivo de esta nota no es repetir la transcripción literal, sino organizar la entrevista en una estructura que se pueda estudiar. La primera línea es la técnica de IA: la ventana de contexto humana, las mejores prácticas de DeepSeek, VLA, el modelo del mundo, Agent OS, Action y la alineación. La segunda línea es la estrategia de la empresa: escala, necesidades de los usuarios, productos tecnológicos, capacidades organizacionales y por qué Li Auto (理想) incluye el automóvil, los robots y el sistema operativo en su visión de los terminales de la era de la AGI. La tercera línea trata de las personas y la organización: la energía, las relaciones íntimas, el cerebro común, el corazón común y la definición de «sabiduría» de Li Xiang.

\begin{figure}[H]
\centering
\includegraphics[width=0.96\textwidth]{ceo-moe-interview-map.png}
\caption{Mapa de la entrevista al gran modelo CEO: Li Xiang tratado como una arquitectura MoE, en la que se invocan conjuntamente la técnica, la estrategia, la organización y las personas. Diagrama conceptual propio, elaborado a partir de la conversación en 00:01:56--00:02:20 y 00:36:18--00:36:25.}
\end{figure}

\begin{knowledgebox}{Lectura de la figura: por qué MoE organiza todo el episodio}
La intuición de MoE es que «distintos expertos se encargan de distintas capacidades». La entrevista usa esta metáfora para dividir las respuestas de Li Xiang en varios canales de expertos: el experto técnico explica VLA y el modelo del mundo, el experto en estrategia explica la escala y la ruta de plataforma, el experto en organización explica la densidad de talento y la metodología, y al final se vuelve a la energía, las relaciones y la sabiduría de una persona. Al leer este episodio, lo importante no es memorizar todas las opiniones, sino observar cómo estos canales de expertos se invocan entre sí.
\end{knowledgebox}

\begin{importantbox}{Tesis central de este episodio}
En la era de la IA, una empresa no se limita a conectar un modelo a un producto: tiene que reescribir al mismo tiempo la forma del producto, la estructura organizacional, el ciclo cerrado de acción y las relaciones humanas. El automóvil, los robots, Agent OS, el modelo del mundo y las relaciones íntimas parecen temas dispersos, pero en realidad responden a la misma pregunta: cómo entra la inteligencia en el mundo real y genera valor de forma continua.
\end{importantbox}

\begin{knowledgebox}{Términos clave: conceptos centrales de este episodio}
\begin{tabularx}{\textwidth}{p{0.20\textwidth}p{0.30\textwidth}X}
\toprule
Término & Problema que resuelve & Significado en este episodio \\
\midrule
MoE & Repartir las capacidades del modelo entre varios expertos & Designa la arquitectura de modelos como DeepSeek y también sirve de metáfora de las capacidades de múltiples expertos de un CEO. \\
VLA & Pasar de la visión y el lenguaje a la acción & Vision-Language-Action, la ruta técnica central con la que Li Auto aborda la conducción autónoma y los robots. \\
World Model & Predecir los cambios del entorno & En escenarios de tráfico se usa para evaluar, generar datos de entrenamiento y podría convertirse en el sistema de operación de la conducción autónoma. \\
Agent OS & Expresión alternativa al Agent general & Li Xiang considera que a corto plazo no habrá un único Agent general, pero sí un sistema operativo en el que funcionen Agents especializados. \\
Action & Que la IA cambie realmente el mundo & Cuando la IA invoca herramientas, activos, vehículos y computadoras, la alineación, los permisos y la responsabilidad adquieren más peso. \\
\bottomrule
\end{tabularx}
\end{knowledgebox}

\subsection{Resumen de la sección}

La dificultad de EP118 está en su enorme amplitud. Al desglosarlo con la metáfora MoE, la línea principal se vuelve mucho más clara: el experto técnico responde «cómo entra el modelo en el mundo físico», el experto en estrategia responde «cómo se convierte la empresa en la plataforma de la siguiente generación», y la organización y las personas responden «quién soporta este cambio».
